% !Mode:: "TeX:UTF-8"
%  Authors: 张井   Jing Zhang: prayever@gmail.com     天津大学2010级管理与经济学部信息管理与信息系统专业硕士生
%           余蓝涛 Lantao Yu: lantaoyu1991@gmail.com  天津大学2008级精密仪器与光电子工程学院测控技术与仪器专业本科生


%%%%%%%%%%%%%%%%% Fonts Definition and Basics %%%%%%%%%%%%%%%%%
%\newcommand{\song}{\CJKfamily{song}}    % 宋体
%\newcommand{\fs}{\CJKfamily{fs}}        % 仿宋体
%\newcommand{\kai}{\CJKfamily{kai}}      % 楷体
%\newcommand{\hei}{\CJKfamily{hei}}      % 黑体
%\newcommand{\li}{\CJKfamily{li}}        % 隶书
\newcommand{\song}{\songti}    % 宋体
\newcommand{\fs}{\fangsong}        % 仿宋体
\newcommand{\kai}{\kaishu}      % 楷体
\newcommand{\hei}{\heiti}      % 黑体
\newcommand{\li}{\lishu}        % 隶书
\newcommand{\yihao}{\fontsize{26pt}{26pt}\selectfont}       % 一号, 单倍行距
\newcommand{\xiaoyi}{\fontsize{24pt}{24pt}\selectfont}      % 小一, 单倍行距
\newcommand{\erhao}{\fontsize{22pt}{1.25\baselineskip}\selectfont}       % 二号, 1.25倍行距
\newcommand{\xiaoer}{\fontsize{18pt}{18pt}\selectfont}      % 小二, 单倍行距
\newcommand{\sanhao}{\fontsize{16pt}{16pt}\selectfont}      % 三号, 单倍行距
\newcommand{\sanhaoup}{\fontsize{16pt}{20pt}\selectfont}    % 三号, 1.25倍行距
\newcommand{\xiaosan}{\fontsize{15pt}{15pt}\selectfont}     % 小三, 单倍行距
\newcommand{\sihao}{\fontsize{14pt}{14pt}\selectfont}       % 四号, 单倍行距
\newcommand{\xiaosi}{\fontsize{12pt}{12pt}\selectfont}      % 小四, 单倍行距
\newcommand{\xiaosiup}{\fontsize{12pt}{15pt}\selectfont}      % 小四, 单倍行距
\newcommand{\wuhao}{\fontsize{10.5pt}{10.5pt}\selectfont}   % 五号, 单倍行距
\newcommand{\xiaowu}{\fontsize{9pt}{9pt}\selectfont}        % 小五, 单倍行距

%\CJKtilde  % 重新定义了波浪符~的意义
% JUST DON'T USE CJK
% 使用 ctexbook 之后已无必要
\newcommand\prechaptername{第}
\newcommand\postchaptername{章}

\punctstyle{hangmobanjiao}             % 调整中文字符的表示，行内占一个字符宽度，行尾占半个字符宽度

% 调整罗列环境的布局
\setitemize{leftmargin=3em,itemsep=0em,partopsep=0em,parsep=0em,topsep=-0em}
\setenumerate{leftmargin=3em,itemsep=0em,partopsep=0em,parsep=0em,topsep=0em}

% 避免宏包 hyperref 和 arydshln 不兼容带来的目录链接失效的问题。
\def\temp{\relax}
\let\temp\addcontentsline
\gdef\addcontentsline{\phantomsection\temp}

% 自定义项目列表标签及格式 \begin{publist} 列表项 \end{publist}
\newcounter{pubctr} %自定义新计数器
\newenvironment{publist}{%%%%%定义新环境

  \begin{list}{[\arabic{pubctr}]} %%标签格式
    {
      \usecounter{pubctr}
      \setlength{\leftmargin}{2.5em}   % 左边界 \leftmargin =\itemindent + \labelwidth + \labelsep
      \setlength{\itemindent}{0em}     % 标号缩进量
      \setlength{\labelsep}{1em}       % 标号和列表项之间的距离,默认0.5em
      \setlength{\rightmargin}{0em}    % 右边界
      \setlength{\topsep}{0ex}         % 列表到上下文的垂直距离
      \setlength{\parsep}{0ex}         % 段落间距
      \setlength{\itemsep}{0ex}        % 标签间距
      \setlength{\listparindent}{0pt}  % 段落缩进量
  }}
{\end{list}}

\makeatletter
\renewcommand\normalsize{
  \@setfontsize\normalsize{12pt}{12pt} % 小四对应 12 pt
  \setlength\abovedisplayskip{4pt}
  \setlength\abovedisplayshortskip{4pt}
  \setlength\belowdisplayskip{\abovedisplayskip}
  \setlength\belowdisplayshortskip{\abovedisplayshortskip}
  \let\@listi\@listI
}

\def\defaultfont{\renewcommand{\baselinestretch}{1.63}\normalsize\selectfont} % 设置行距
%
\renewcommand{\CJKglue}{\hskip -0.1 pt plus 0.08\baselineskip} % 控制字间距，使每行 34 个汉字
\makeatother

%%%%%%%%%%%%% Contents %%%%%%%%%%%%%%%%%
\renewcommand{\contentsname}{\centering \xiaosan \heiti 目\quad 录}
\setcounter{tocdepth}{2} % 控制目录深度

\titlecontents{chapter}[2em]
  {\vspace{.25\baselineskip}\xiaosi\songti}
  {\heiti\thecontentslabel\quad \songti}
  {}
  {\hspace{.3em}\titlerule*{.}\contentspage}

\titlecontents{section}[3em]
  {\vspace{.25\baselineskip}\xiaosi\songti}
  {\heiti\thecontentslabel\quad \songti}
  {}
  {\hspace{.5em}\titlerule*{.}\contentspage}
%
\titlecontents{subsection}[4em]
  {\vspace{.25\baselineskip}\xiaosi\songti}
  {\heiti\thecontentslabel\quad \songti}
  {}
  {\hspace{.5em}\titlerule*{.}\contentspage}

%%%%%%%%%% Chapter and Section %%%%%%%%%%%%%
\setcounter{secnumdepth}{4}
\setlength{\parindent}{1em}
\renewcommand{\chaptername}{\prechaptername\CJKnumber{\thechapter}\postchaptername}

% Chapter 第一章 绪论 [xiaosan, heiti]
\titleformat{\chapter}[block]
  {\centering\xiaosan\heiti} % format
  {} % label
  {0pt} % sep
  {} % before-code
\titlespacing{\chapter}
  {0pt} % 左边距
  {0pt} % 上面的空白
  {20pt} % 下面的空白

% Section
\titleformat{\section}{\sihao\heiti}{\thesection}{1em}{}
\titlespacing{\section}{0pt}{18pt}{18pt}

% Subsection
\titleformat{\subsection}{\sihao\heiti}{\thesubsection}{1em}{}
\titlespacing{\subsection}{0pt}{12pt}{12pt}

% SubsubSection
\titleformat{\subsubsection}{\sihao\heiti}{\thesubsubsection}{1em}{}
\titlespacing{\subsubsection}{0pt}{10pt}{10pt}

%%%%%%%%%% Table, Figure and Equation %%%%%%%%%%%%%%%%%
\renewcommand{\tablename}{表}                                     % 插表题头
\renewcommand{\figurename}{图}                                    % 插图题头
\renewcommand{\thefigure}{\arabic{chapter}-\arabic{figure}}       % 使图编号为 7-1 的格式 %\protect{~}
\renewcommand{\thesubfigure}{\alph{subfigure}}                   % 使子图编号为 subfig 自动就是 (a) 的格式 这里保留只是作为参考
\renewcommand{\thesubtable}{(\alph{subtable})}                    % 使子表编号为 (a) 的格式
\renewcommand{\thetable}{\arabic{chapter}-\arabic{table}}         % 使表编号为 7-1 的格式
\renewcommand{\theequation}{\arabic{chapter}-\arabic{equation}}   % 使公式编号为 7-1 的格式
\newcommand{\ud}{\mathrm{d}}

%%%%%% 定制浮动图形和表格标题样式 %%%%%%
\makeatletter
\long\def\@makecaption#1#2{
  \vskip\abovecaptionskip
  \sbox\@tempboxa{\centering\wuhao\song{#1\qquad #2} }
  \ifdim \wd\@tempboxa >\hsize
    \centering\wuhao\song{#1\qquad #2} \par
  \else
    \global \@minipagefalse
    \hb@xt@\hsize{\hfil\box\@tempboxa\hfil}
  \fi
\vskip\belowcaptionskip}
\makeatother
\captiondelim{~~~~} %用来控制longtable表头分隔符

%%%%%%%%%% Theorem Environment %%%%%%%%%%%%%%%%%
\theoremstyle{plain}
\theorembodyfont{\song\rmfamily}
\theoremheaderfont{\hei\rmfamily}
\newtheorem{theorem}{定理~}[chapter]
\newtheorem{lemma}{引理~}[chapter]
\newtheorem{axiom}{公理~}[chapter]
\newtheorem{proposition}{命题~}[chapter]
\newtheorem{prop}{性质~}[chapter]
\newtheorem{corollary}{推论~}[chapter]
\newtheorem{definition}{定义~}[chapter]
\newtheorem{conjecture}{猜想~}[chapter]
\newtheorem{example}{例~}[chapter]
\newtheorem{remark}{注~}[chapter]
%\newtheorem{algorithm}{算法~}[chapter]
\newenvironment{proof}{\noindent{\hei 证明：}}{\hfill $ \square $ \vskip 4mm}
\theoremsymbol{$\square$}

%%%%%%%%%% Page: number, header and footer  %%%%%%%%%%%%%%%%%

%\frontmatter 或 \pagenumbering{roman}
%\mainmatter 或 \pagenumbering{arabic}

%%%%%%%%%%% Code: Listings from MCM Template %%%%%%%%%%%%

\definecolor{grey}{rgb}{0.8,0.8,0.8}
\definecolor{darkgreen}{rgb}{0,0.3,0}
\definecolor{darkblue}{rgb}{0,0,0.3}
\def\lstbasicfont{\fontfamily{pcr}\selectfont\footnotesize}
\lstset{%
  % indexing
  % numbers=left,
  % numberstyle=\small,%
  % character display
  showstringspaces=false,
  showspaces=false,%
  tabsize=4,%
  % style
  frame=lines,%
  basicstyle={\footnotesize\lstbasicfont},%
  keywordstyle=\color{darkblue}\bfseries,%
  identifierstyle=,%
  commentstyle=\color{darkgreen},%\itshape,%
  stringstyle=\color{black}%
}
\lstloadlanguages{C,C++,Java,Matlab,Mathematica,Python}

%%%%%%%%%%%% References %%%%%%%%%%%%%%%%%
\renewcommand{\bibname}{参考文献}

% 重定义参考文献样式，来自thu
\makeatletter
\renewenvironment{thebibliography}[1]{
  \addcontentsline{toc}{chapter}{参考文献}
  \chapter*{\bibname}
  \wuhao
  \list{\@biblabel{\@arabic\c@enumiv}}
  {\renewcommand{\makelabel}[1]{##1\hfill}
    \settowidth\labelwidth{0 cm}
    \setlength{\labelsep}{0pt}
    \setlength{\itemindent}{0pt}
    \setlength{\leftmargin}{\labelwidth+\labelsep}
    \addtolength{\itemsep}{-0.7em}
    \usecounter{enumiv}
    \let\p@enumiv\@empty
  \renewcommand\theenumiv{\@arabic\c@enumiv}}
  \sloppy\frenchspacing
  \clubpenalty4000
  \@clubpenalty \clubpenalty
  \widowpenalty4000
  \interlinepenalty4000
\sfcode`\.\@m}
{\def\@noitemerr
  {\@latex@warning{Empty `thebibliography' environment}}
\endlist\frenchspacing}
\makeatother

\addtolength{\bibsep}{-0.5em}     % 缩小参考文献间的垂直间距
\setlength{\bibhang}{2em}         % 每个条目自第二行起缩进的距离

% 参考文献引用作为上标出现
% \newcommand{\citeup}[1]{\textsuperscript{\cite{#1}}}
\makeatletter
\def\@cite#1#2{\textsuperscript{[{#1\if@tempswa , #2\fi}]}}
\makeatother
%% 引用格式
\bibpunct{[}{]}{,}{s}{}{,}

%%%%%%%%%%%% Cover %%%%%%%%%%%%%%%%%
% 封面、摘要、版权、致谢格式定义
\makeatletter
\def\ctitle#1{\def\@ctitle{#1}}\def\@ctitle{}
\def\cdegree#1{\def\@cdegree{#1}}\def\@cdegree{}
\def\caffil#1{\def\@caffil{#1}}\def\@caffil{}
\def\csubject#1{\def\@csubject{#1}}\def\@csubject{}
\def\cgrade#1{\def\@cgrade{#1}}\def\@cgrade{}
\def\cauthor#1{\def\@cauthor{#1}}\def\@cauthor{}
\def\cnumber#1{\def\@cnumber{#1}}\def\@cnumber{}
\def\csupervisor#1{\def\@csupervisor{#1}}\def\@csupervisor{}
\def\crank#1{\def\@crank{#1}}\def\@crank{}
\def\cdeclaration#1{\def\@cdeclaration{#1}}\def\@cdeclaration{}
\def\cdate#1{\def\@cdate{#1}}\def\@cdate{}
\long\def\cabstract#1{\long\def\@cabstract{#1}}\long\def\@cabstract{}
\long\def\eabstract#1{\long\def\@eabstract{#1}}\long\def\@eabstract{}
\def\ckeywords#1{\def\@ckeywords{#1}}\def\@ckeywords{}
\def\ekeywords#1{\def\@ekeywords{#1}}\def\@ekeywords{}
\def\cheading#1{\def\@cheading{#1}}\def\@cheading{}


\pagestyle{fancy}
\fancyhf{}
\fancyhead[C]{\song\wuhao \@cheading}  % 页眉显示天津大学 20XX 届本科生毕业论文
\fancyfoot[C]{\song\xiaowu ~\thepage~}
\newlength{\@title@width}

% 定义封面
\def\makecover{
  %\cleardoublepage%
  \phantomsection
  \pdfbookmark[-1]{\@ctitle}{ctitle}

  \begin{titlepage}
    \begin{center}

      \begin{figure}[h]
        \centering
        \includegraphics[width=0.4\textwidth]{figures/tjuname.eps}
      \end{figure}
      \hei\erhao{\textbf{本科生毕业论文}}
      \vspace*{60pt}

      \begin{figure}[h]
        \centering
        \includegraphics[width=0.4\textwidth]{figures/tjulogo.eps}
      \end{figure}

      \vspace*{10pt}

      \heiti\sanhaoup\textbf{题目：\@ctitle}

      \vspace{30pt}

      \renewcommand\arraystretch{1.5}
      \setlength{\@title@width}{5cm}
      {\sanhaoup\song{\bf{
            \begin{tabular}{lc}
              学\qquad 院&  \underline{\makebox[\@title@width][c]{\@caffil}} \\
              专\qquad 业 &  \underline{\makebox[\@title@width][c]{\@csubject}} \\
              年\qquad 级  &  \underline{\makebox[\@title@width][c]{\@cgrade}}\\
              姓\qquad 名 &  \underline{\makebox[\@title@width][c]{\@cauthor}} \\
              学\qquad 号 &  \underline{\makebox[\@title@width][c]{\@cnumber}} \\
              指导教师 &  \underline{\makebox[\@title@width][c]{\@csupervisor}} \\
        \end{tabular}}}
      }
    \end{center}
  \end{titlepage}

  %%%%%%%%%%%%%%%%%%%   独创性声明  %%%%%%%%%%%%%%%%%%%%%%%
  \vspace{27.5pt}
  \markboth{独创性声明}{独创性声明}
  \pdfbookmark[0]{独创性声明}{cdeclaration}
  \chapter*{\textbf{\centering\erhao\songti 独创性声明}}
  \vspace{27.5pt}
  \sanhaoup % 三号字体 1.25 倍行距
  \@cdeclaration 
  \vspace{40pt} % 两行 三号字体 1.25 倍行距

  \begin{center}
    论文作者签名:
  \end{center}

  \hfill  年 \hspace{1cm} 月 \hspace{1cm} 日

  \vspace{40pt} % 两行 三号字体 1.25 倍行距
  本人声明：本毕业设计（论文）是本人指导学生完成的研究成果，已经审阅过论文的全部内容。

  \vspace{20pt} % 两行 三号字体 1.25 倍行距
  \begin{center}
    论文指导教师签名:
  \end{center}

  \hfill  年 \hspace{1cm} 月 \hspace{1cm} 日

  \thispagestyle{empty}

  \pagenumbering{Roman} % 页脚页码格式 - 大写罗马字母
  %%%%%%%%%%%%%%%%%%%   Abstract and Keywords  %%%%%%%%%%%%%%%%%%%%%%%
  % \clearpage
  \markboth{摘~要}{摘~要}
  \pdfbookmark[0]{摘~~要}{cabstract}
  \chapter*{\textbf{\centering\erhao\songti 摘\quad 要}}
  \setcounter{page}{1}                                 % 单独从 1 开始编页码

  \fontsize{12pt}{20pt}\selectfont % 小四号，20pt 固定行距
  \@cabstract

  \vspace{\baselineskip}

  \noindent\textbf{\songti\sihao 关键词:} \@ckeywords

  % \thispagestyle{empty} % 去掉中文摘要页眉页脚

  %%%%%%%%%%%%%%%%%%%   English Abstract  %%%%%%%%%%%%%%%%%%%%%%%%%%%%%%
  % \clearpage
  %\phantomsection
  \markboth{ABSTRACT}{ABSTRACT}
  \pdfbookmark[0]{ABSTRACT}{eabstract}
  \chapter*{\centering\erhao{\bf{ABSTRACT}}}

  \fontsize{12pt}{15pt}\selectfont % 小四号，20pt 固定行距

  \@eabstract
  \vspace{\baselineskip}

  % \hangafter=1\hangindent=60pt
  \noindent\textbf{\sihao KEYWORDS:} \@ekeywords

  % \thispagestyle{empty} % 去掉英文摘要页眉页脚
}

\makeatother

%%%%%%%%%%%%%%%%%%%% 嵌入单张图片宏 %%%%%%%%%%%%%%%%%%%%
\newcommand{\figuremacro}[5]{
  \begin{figure}[#1]
    \centering
    \includegraphics[width=#5\columnwidth]{#2}
    \caption[#3]{\textbf{#3}#4}
    \label{fig:#2}
  \end{figure}
}
%% \figuremacro{ht}{Method1}{Illustration 1}{ - The Colored Red and Light blue patterns are moved from top to bottom and left to right.}{1.0}

%%%%%%%%%%%%%%%%%%%%%%%%%%% Hack %%%%%%%%%%%%%%%%%%%%
\makeatletter
\patchcmd{\src@startchapter}{\if@openright\cleardoublepage\else\clearpage\fi}{}{}{}
\makeatother
